\documentclass[11pt]{article}
\usepackage{mathblatt}
\makeatletter
\setlength{\mbpfbe}{0mm}% keine Punkte (Bauregel 6.4)
\makeatother
% Rechenraum als Karo (Bauregel 6.9): n Rechenschritte = n mal 10 mm
\newcommand{\karo}[1]{\par\smallskip\noindent\begin{tikzpicture}
  \clip (0,0) rectangle ({\linewidth-1mm},{#1*10mm});
  \draw[black!16,line width=0.3pt,step=5mm] (0,0) grid ({\linewidth},{#1*10mm});
  \end{tikzpicture}\par}
\newcommand{\linie}[1]{\,{\color{black!45}\rule[-1pt]{#1}{0.4pt}}\,}
\newcommand{\pkt}{\textperiodcentered{}}
\begin{document}

\pfheftstil
\pflaufkopf[P10 \pkt{} mehr]{Lineare Funktionen \pkt{} Lösungen}{PA10-6}\pfseitevon
\pfheftkopf{Lineare Funktionen}{Lösungen \pkt{} mehr}
\begin{pfloesungb}{}
\lzb{1.}{}{y = $-$6}{y = $-$6}
\end{pfloesungb}
\begin{pfloesungb}{}
\lzb{2.}{}{0 min: 36 cm; 80 min: 12 cm}{0 min: 36 cm; 80 min: 12 cm}
\end{pfloesungb}
\begin{pfloesungb}{}
\lzb{3.}{}{120 min}{120 min}
\end{pfloesungb}
\begin{pfloesungb}{}
\lzb{4.}{}{Fitnesscenter $\to$ y = 4x + 15}{Fahrpreis $\to$ y = 0,30x + 4; Fitnesscenter $\to$ y = 4x + 15}
\end{pfloesungb}
\begin{pfloesungb}{}
\lzb{5.}{}{II beginnt bei 150 € (Grundpreis bei 0 km) und steigt flacher (1,20 €/km)}{I Reiselust, II Sonnenschein; II beginnt bei 150 € (Grundpreis bei 0 km) und steigt flacher (1,20 €/km)}
\end{pfloesungb}
\begin{pfloesungb}{}
\lzb{6.}{}{Gerade durch (0|$-$1) und (1|1)}{Gerade durch (0|$-$1) und (1|1)}
\end{pfloesungb}
\begin{pfloesungb}{}
\lzb{7.}{}{Aussage 4 falsch}{Aussage 1 falsch; Aussage 2 wahr; Aussage 3 wahr; Aussage 4 falsch}
\end{pfloesungb}
\begin{pfloesungb}{}
\lzb{8.}{}{36: Anfangshöhe der Kerze in cm}{y: Höhe der Kerze in cm; x: Zeit in Minuten seit dem Anzünden; 36: Anfangshöhe der Kerze in cm}
\end{pfloesungb}
\begin{pfloesungb}{}
\lzb{9.}{}{Reiselust billiger: Sonnenschein 414 € gegen Reiselust 396; nein: bei 280 km Reiselust 504 €, Sonnenschein 486}{Reiselust billiger: Sonnenschein 414 € gegen Reiselust 396; nein: bei 280 km Reiselust 504 €, Sonnenschein 486}
\end{pfloesungb}
\begin{pfloesungb}{}
\lzb{10.}{}{x = $-$3}{x = $-$3}
\end{pfloesungb}
\begin{pfloesungb}{}
\lzb{11.}{}{y = 1,2x + 150 (x in km, y in €)}{y = 1,2x + 150 (x in km, y in €)}
\end{pfloesungb}
\begin{pfloesungb}{}
\lzb{12.}{}{f}{f}
\end{pfloesungb}
\begin{pfloesungb}{}
\lzb{13.}{}{Gerade durch K und L (y-Achsenabschnitt 2, Nullstelle $-$3); n = $\tfrac{2}{3}$x + 2}{Gerade durch K und L (y-Achsenabschnitt 2, Nullstelle $-$3); m = 6 : 9 = $\tfrac{2}{3}$; n = 4 $-$ $\tfrac{2}{3}$ · 3 = 2, also y = $\tfrac{2}{3}$x + 2}
\end{pfloesungb}
\begin{pfloesungb}{}
\lzb{14.}{}{h(t) = $-$6t + 1650}{h(t) = $-$6t + 1650}
\end{pfloesungb}
\end{document}